\documentclass[10pt,letterpaper]{article}
\usepackage{cornell-notes}

\title{Chapter 7: Oscillators and Timers}
\author{}
\date{}
\setDocTitle{Chapter 7: Oscillators and Timers}
\setDocSubtitle{Electronics}
\setDocOwner{Electronics Cornell Notes}
\setDocDate{\today}
\setCornellCollection{Electronics Cornell Notes}
\setCornellUnitType{Chapter}
\setCornellUnitNumber{7}
\setCornellUnitTitle{Oscillators and Timers}
\hypersetup{pdftitle={Chapter 7: Oscillators and Timers},pdfsubject={Cornell notes},pdfauthor={}}
\begin{document}
\maketitle

\begin{CornellOverviewBox}{Concept}
Oscillators create periodic signals without an external periodic input; timers create controlled intervals or pulse widths from events. Designs range from RC relaxation and 555 circuits to sine oscillators, crystals, DDS, PLL synthesis, quadrature sources, monostables, and counter-based timing. Frequency accuracy, startup, amplitude control, jitter, retrigger behavior, and recovery are as important as the nominal timing equation.
\end{CornellOverviewBox}
\begin{CornellOverviewBox}{Concept}
\textbf{Prerequisites.} RC charging, comparator hysteresis, op-amp feedback, resonance, logic levels, and basic probability/noise concepts.

\textbf{After studying this chapter, you should be able to:}
\begin{itemize}
  \item Explain startup and steady-state conditions for relaxation and sinusoidal oscillators.
  \item Calculate timer periods while accounting for thresholds, tolerance, leakage, and propagation delay.
  \item Compare crystal, TCXO/OCXO, DDS, PLL, and counter-based timing for accuracy and jitter.
  \item Design monostable and pulse-conditioning circuits with explicit trigger and retrigger behavior.
\end{itemize}\textbf{Central design questions.}
\begin{itemize}
  \item What determines nominal frequency, long-term accuracy, and short-term jitter?
  \item How does amplitude settle without clipping or dying out?
  \item What happens for slow, noisy, repeated, or missing triggers and during power-up?
\end{itemize}\end{CornellOverviewBox}
\section{Cornell Notes}
\begin{CornellNotesTable}
\CornellNoteRow{7.1.1: Oscillation condition}{A sinusoidal loop starts when noise experiences loop gain above unity with phase near $0^{\circ}$ modulo $360^{\circ}$. In steady state, nonlinear amplitude control reduces effective loop gain to unity. The Barkhausen statement is a useful boundary description, not a complete startup guarantee; pole locations, startup margin, limiting mechanism, and noise determine actual behavior.}
\CornellNoteRow{7.1.2: Relaxation oscillators}{A capacitor charges between thresholds while a comparator or Schmitt trigger reverses current or polarity. Frequency depends on $RC$ and threshold ratios. Comparator delay, output saturation, input common-mode limits, hysteresis tolerance, leakage, and slew create error. Keep threshold reference quiet and ensure the capacitor current dominates leakage.}
\CornellNoteRow{7.1.3--7.1.4: 555 and timer ICs}{A classic timer compares a capacitor with internal fractional-supply thresholds and controls discharge/output stages. Astable, monostable, and bistable modes are convenient, but control-pin noise, discharge saturation, load-current supply bounce, timing-capacitor leakage, and bipolar input currents limit precision. CMOS versions reduce current and leakage but have different drive and speed.}
\CornellNoteRow{7.1.5: Sine oscillators}{RC bridges, phase-shift networks, and LC resonators provide frequency-selective positive feedback. Amplitude control may use lamps, diodes, FETs, AGC, or controlled limiting. Hard clipping creates harmonics; slow AGC can cause amplitude bounce. Component tolerance and amplifier phase shift move frequency. Design startup gain slightly above unity, then stabilize amplitude smoothly.}
\CornellNoteRow{7.1.6--7.1.7: Crystal stability}{A quartz crystal behaves like a high-$Q$ motional resonator plus shunt capacitance. Load capacitance, drive level, oscillator topology, aging, temperature, and board contamination affect frequency. TCXOs compensate temperature; OCXOs control resonator temperature for better stability at higher power and warm-up time. Never overdrive or assume a crystal's labeled frequency is topology-independent.}
\CornellNoteRow{7.1.8: DDS and PLL synthesis}{DDS uses a phase accumulator, lookup/conversion, and clock to produce fine frequency steps; spurs, phase truncation, DAC images, and reference-clock quality matter. A PLL compares phase, filters error, and controls an oscillator; division multiplies reference frequency and also transforms phase-noise relationships. Loop bandwidth trades reference tracking, VCO noise suppression, lock time, and spur rejection.}
\CornellNoteRow{7.1.9: Quadrature}{Quadrature sources provide signals separated by $90^{\circ}$. Polyphase/allpass networks, state-variable oscillators, digital division, or DDS can generate them. Evaluate amplitude balance, phase error versus frequency, startup state, and common clocking. A divide-by-four digital method can produce accurate phase but quantized waveform and frequency constraints.}
\CornellNoteRow{7.1.10: Jitter}{Jitter is time variation of threshold crossings; phase noise is its frequency-domain description. Noise near a waveform's crossing converts to timing uncertainty roughly as $\sigma_t\approx\sigma_v/(dv/dt)$. Faster slew reduces threshold-noise sensitivity, while supply noise, comparator noise, reference noise, and oscillator phase noise add. Specify observation bandwidth and whether jitter is period, cycle-to-cycle, or accumulated.}
\CornellNoteRow{7.2.1--7.2.3: Monostables}{A trigger initiates one output interval set by an analog ramp or digital count. Define trigger polarity, minimum pulse, retriggerable/nonretriggerable behavior, timeout extension, output polarity, and response to a trigger arriving during timeout. Analog monostables inherit RC/leakage/threshold errors; dedicated ICs can have minimum timing capacitance and power-up quirks.}
\CornellNoteRow{7.2.4: Counter timing}{A stable clock and counter create intervals $T=N/f_{clk}$ with high repeatability and programmable range. Quantization is one clock period, while synchronization and metastability matter for asynchronous triggers. Counter width limits maximum interval; clock tolerance sets absolute accuracy. Use safe clock-domain crossing and define reset state.}
\end{CornellNotesTable}
\begin{CornellSummaryBox}{Summary}
Timing circuits combine a frequency-selective or integrating element with gain, thresholds, and nonlinear state control. Nominal $RC$, $LC$, or count equations are only the beginning. Accuracy and spectral purity depend on reference quality, thresholds, delay, leakage, amplitude control, jitter, triggering semantics, power-up state, and the measurement interval.
\end{CornellSummaryBox}
\section{Worked Examples}
\begin{CornellExamBox}{Schmitt relaxation estimate}
A symmetric inverting relaxation oscillator charges through $R=68\,\mathrm{k\Omega}$ into $C=10\,\mathrm{nF}$ between thresholds $\pm\beta V_o$ with $\beta=0.4$. The period is $T=2RC\ln[(1+\beta)/(1-\beta)]=1.15\,\mathrm{ms}$, so $f\approx868\,\mathrm{Hz}$. Output swing and threshold ratios must use actual loaded values.
\end{CornellExamBox}
\begin{CornellExamBox}{Counter interval}
A $12\,\mathrm{MHz}$ clock divided by $N=36000$ gives $T=3.000\,\mathrm{ms}$. With clock tolerance $\pm25\,\mathrm{ppm}$, frequency-derived interval error is $\pm75\,\mathrm{ns}$, plus trigger synchronization quantization up to one clock period ($83.3\,\mathrm{ns}$).
\end{CornellExamBox}
\begin{CornellExamBox}{Noise-to-jitter}
A crossing has slope $0.8\,\mathrm{V/\mu s}$ and rms threshold-referred noise $1.6\,\mathrm{mV}$. The first timing estimate is $\sigma_t=1.6\,\mathrm{mV}/(0.8\,\mathrm{V/\mu s})=2.0\,\mathrm{ns}$. This excludes correlated oscillator phase noise and bandwidth-definition effects.
\end{CornellExamBox}
\section{Design and Troubleshooting}
\begin{CornellExamBox}{Design and Troubleshooting}
\begin{enumerate}
  \item Specify waveform, frequency/range, accuracy, phase noise/jitter, startup, amplitude, load, and power.
  \item Choose RC, LC, crystal, DDS, PLL, or counter architecture from stability and tuning needs.
  \item Budget reference error, thresholds, delay, leakage, noise, temperature, aging, and quantization.
  \item Prototype and measure startup, frequency versus supply/temperature, spectrum, jitter, trigger corner cases, and power-cycle state.
\end{enumerate}\end{CornellExamBox}
\begin{CornellExamBox}{Symptoms, likely causes, and checks}
\begin{itemize}
  \item No startup: loop gain too small, wrong phase, resonator loading, amplifier range violation, or excessive damping.
  \item Frequency wrong: actual thresholds/output swing, load capacitance, parasitics, clock setting, or unit error.
  \item High distortion: hard limiting, insufficient amplifier bandwidth/slew, or poor AGC.
  \item False triggers: slow/noisy edge, inadequate hysteresis, ground bounce, or asynchronous sampling.
  \item Large jitter: low crossing slew, noisy supply/reference, comparator noise, PLL bandwidth choice, or poor clock source.
\end{itemize}\end{CornellExamBox}
\begin{CornellWarningBox}{Safety and Limitations}
Oscillator outputs can drive resonant or power stages to destructive voltage/current even from a low-energy control circuit. Limit drive during startup tests, rate capacitors and inductors for peak stress, control EMI, and avoid probing high-voltage resonant nodes with inappropriate grounds or probe ratings.
\end{CornellWarningBox}
\section{Key-Equation Sheet}
\begin{CornellExamBox}{Key Equations}
\begin{itemize}
  \item Relaxation timing generally has $T=RC$ times a logarithm of threshold ratios; derive it from the actual charge endpoints.
  \item Ideal resonance: $f_0=1/(2\pi\sqrt{LC})$.
  \item Counter timing: $T=N/f_{clk}$ and resolution is one clock period.
  \item DDS output: $f_{out}=Kf_{clk}/2^M$ for tuning word $K$ and $M$-bit phase accumulator.
  \item PLL multiplication: $f_{out}=(N/R)f_{ref}$ for divider conventions $N$ and $R$.
  \item Threshold-noise timing estimate: $\sigma_t\approx\sigma_v/(dv/dt)$ at the crossing.
\end{itemize}\end{CornellExamBox}
\section{Study and Review}
\subsection{Design checklist}
\begin{itemize}
  \item Startup and amplitude-control margin are defined.
  \item Reference, $RC/LC$, threshold, delay, temperature, aging, and jitter budgets pass.
  \item Trigger polarity, retrigger, reset, and power-up behavior are explicit.
  \item Load, output spectrum, EMI, and supply-current transients are acceptable.
  \item Frequency and jitter measurements use sufficient observation time and bandwidth.
\end{itemize}\subsection{Common misconceptions}
\begin{itemize}
  \item Loop gain exactly one does not guarantee startup.
  \item A 555 timing equation does not include every leakage and threshold error.
  \item A crystal is not an ideal fixed-frequency component independent of load.
  \item Fine DDS resolution does not imply low spurs.
  \item Jitter has multiple definitions and must include bandwidth/interval.
\end{itemize}\subsection{Active-recall questions}
\begin{enumerate}
  \item What two loop conditions support sinusoidal oscillation?
  \item How does amplitude stabilize?
  \item Which terms set relaxation frequency?
  \item Why does crystal load capacitance matter?
  \item What does a PLL loop bandwidth trade?
  \item How does DDS frequency tuning work?
  \item What converts voltage noise to jitter?
  \item How do retriggerable and nonretriggerable monostables differ?
  \item What is counter timing quantization?
  \item Which test proves reliable startup?
\end{enumerate}\subsection{Original practice problems}
\begin{enumerate}
  \item A $16\,\mathrm{MHz}$ clock must create $250\,\mu\mathrm{s}$. Find the count.
  \item Find ideal $LC$ resonance for $L=10\,\mu\mathrm{H}$ and $C=220\,\mathrm{pF}$.
  \item A 32-bit DDS clocked at $100\,\mathrm{MHz}$ uses tuning word 42,949,673. Estimate output.
\end{enumerate}\subsection{Answer key / solution outline}
\begin{enumerate}
  \item $N=Tf=4000$ counts; synchronization can add up to a clock of timing uncertainty.
  \item $f_0=3.39\,\mathrm{MHz}$; stray capacitance may be a significant fraction.
  \item $f=Kf_c/2^{32}\approx1.000\,\mathrm{MHz}$; spectral purity still depends on clock/DAC/filtering.
\end{enumerate}\begin{CornellSummaryBox}{Summary}
A timing equation predicts nominal period; a timing design also proves startup, amplitude control, reference stability, jitter, trigger semantics, reset state, load behavior, and environmental drift.
\end{CornellSummaryBox}

\end{document}
